\chapter{Conclusion and Future Perspectives}
The availability of cheap predictive tools helps advance science in a number of ways. First and foremost, good models save costs by narrowing down the number of hypotheses. Synthetic enhancer candidates, for example, can be evaluated in-silico, and only the most promising ones are kept in the reporter cohort. The resulting higher iteration speed and reduced capital cost also reduce the barrier to entry, allowing us to utilise human resources more efficiently. Finally, since ML models value negative and control data, their use incentivizes data sharing and transparency, and potentially the repurposing of existing databases.
\begin{itemize}[noitemsep]
    \item Cell response to drugs
    \item Differentiate cells into what we want, transhumanism
    \item protein binding
\end{itemize}
There are three main areas: the first, and the most logical continuations on this work specifically, is the effect of genetic perturbations on the regulome. As shown in chapter [x], Perturbation screens have already been used as a tool for causative inference. The resulting models can be used to help understand disease ethiology from a genomic or regulomic point of view. Second is cellular response to small molecules in drug design. While medicinal chemists make use of some modeling tools, most work is still performed using brute force screens, intelligent design, and iterative campaigns. Finally, and perhaps the "holy grail" of regulomics, is deterministic modulation of cellular differentiation processes. Using predictive models, we should ultimately be able to develop differentiation processes for various cell types. This will in turn facilitate the development of advanced cell therapies, where we can autologously replace a patient's diseased or aged cells with autologous new ones. [give eaxmples]

Consolidating this knowledge into deterministic understanding will ultimately lead to advanced cell- and gene therapies. CAR-T, the re-engineering of autologous T cells to target cancer [be more specific], has been a resounding success.


so number of eras
\begin{itemize}[noitemsep]
    \item Early era of biology
    \item Early molecular era
    \item Commodified sequencing
    \item Dawn of AI
    \item Future: massive data foundries (cite that paper that shows data still scales super well, maybe bitter pill as well)
\end{itemize}

this phd is really about data generation at large scale
hypothesis: ai has provided an incentive to generate massive datasets in the private sector, because they can now be monetized without direct access! distribution of data

\sjcitep[.]{pampari2025}

% closing statements

AI/ML in bioinformatics, specifically single-cell, is not new.
the "level" of predictions is going higher and higher
first order: low level predictions
second order: feature predictions in individual datasets
now: genome-wide predictions (?)
future: higher level, cross-modal predictions (drug response?)

what do you need for that: scale, fidelity and multiple arrayed conditions

how do you achieve that?
automation, arrayed screens, - liquid handlers (droplets are not good at handling arrayed conditions, except if you can multiplex)
more robust sample prep (remove humans out of the loop)

why is the time right now?
ai can actually use this data
sequencing becoming cheap enough?
vc appetite for these things?

\sjcitep[.]{wierenga2023pylabrobot}